\documentclass{article}
\usepackage{amsmath,amssymb,amsthm}
\usepackage{algorithm}
\usepackage[noend]{algpseudocode}
\usepackage{enumitem}
\usepackage{hyperref}
\newtheorem{theorem}{Theorem}
\newtheorem{lemma}{Lemma}
\newtheorem{assumption}{Assumption}
\newcommand{\E}{\mathbb{E}}
\newcommand{\R}{\mathbb{R}}

\begin{document}

\section*{AdaGrad (Scalar Version)}

\section*{Mathematical Formulation}
We consider the unconstrained minimisation problem  

\[
\min_{w\in\R} \; f(w),
\]

where $f:\R\rightarrow\R$ is convex and $L$‑smooth.  
At iteration $k\ge 1$ the algorithm has the current point $w_k\in\R$ and the
\emph{accumulated squared gradient}
\[
G_k \;:=\; \sum_{i=1}^{k} g_i^{2},\qquad 
g_i \;:=\; \nabla f(w_i).
\]

The learning rate at iteration $k$ is
\[
\eta_k \;:=\; \frac{\alpha}{\sqrt{G_k+\varepsilon}},\qquad 
\alpha>0,\;\varepsilon>0,
\]
and the update rule is
\[
\boxed{w_{k+1}= w_k - \eta_k\, g_k } .
\]

Thus the algorithm maintains a running average of (squared) gradients and
scales the step size inversely proportionally to the root of this average.
The hyper‑parameters are $\alpha$ (initial stepsize) and $\varepsilon$ (regulariser
to avoid division by zero).

\section*{Pseudocode}
\begin{algorithm}[H]
\caption{Scalar AdaGrad}
\begin{algorithmic}[1]
\Require Initial point $w_0\in\R$, stepsize parameter $\alpha>0$, 
        stabiliser $\varepsilon>0$.
\State $G_0 \gets 0$ \Comment{accumulated squared gradient}
\For{$k=0,1,2,\dots,T-1$}
    \State $g_k \gets \nabla f(w_k)$
    \State $G_k \gets G_{k-1}+ g_k^{2}$ \Comment{update accumulator}
    \State $\eta_k \gets \dfrac{\alpha}{\sqrt{G_k+\varepsilon}}$
    \State $w_{k+1} \gets w_k - \eta_k\, g_k$
\EndFor
\State \Return $w_T$ (or any averaged iterate)
\end{algorithmic}
\end{algorithm}

\section*{Dry Run Example}
Minimise $f(w)=(w-3)^2$.
\[
\nabla f(w)=2(w-3),\qquad f\text{ is $L=2$ smooth and convex.}
\]

Take $w_0=0$, $\alpha=0.1$, $\varepsilon=10^{-8}$.
\begin{center}
\begin{tabular}{c|c|c|c|c}
$k$ & $w_k$ & $g_k= \nabla f(w_k)$ & $G_k$ & $w_{k+1}=w_k-\eta_k g_k$\\\hline
0 & $0$ & $2(0-3)=-6$ & $G_0=36$ & $\displaystyle\eta_0=\frac{0.1}{\sqrt{36}}\!=0.016666$\\
  &   &   &   & $w_1=0-0.016666\cdot(-6)=0.1000$\\\hline
1 & $0.1000$ & $2(0.1-3)=-5.8$ & $G_1=36+33.64=69.64$ &
 $\displaystyle\eta_1=\frac{0.1}{\sqrt{69.64}}=0.01198$\\
  &   &   &   & $w_2=0.1-0.01198\cdot(-5.8)=0.1695$\\\hline
2 & $0.1695$ & $2(0.1695-3)=-5.661$ & $G_2=69.64+32.05=101.69$ &
 $\displaystyle\eta_2=\frac{0.1}{\sqrt{101.69}}=0.00993$\\
  &   &   &   & $w_3=0.1695-0.00993\cdot(-5.661)=0.2221$\\\hline
\end{tabular}
\end{center}
Objective values: $f(w_0)=9$, $f(w_1)=8.41$, $f(w_2)=8.04$, $f(w_3)=7.77$.
The step sizes shrink as the accumulated squared gradients increase,
illustrating the AdaGrad adaptive behaviour.

\newpage
\section*{Assumptions}
\begin{assumption}[Convexity]\label{as:convex}
$f:\R\to\R$ is convex, i.e.
\[
f(u)\ge f(v)+\nabla f(v)(u-v),\quad\forall u,v\in\R .
\]
\end{assumption}
\begin{assumption}[Smoothness]\label{as:smooth}
$f$ is $L$‑smooth: for all $u,v\in\R$,
\[
|\,\nabla f(u)-\nabla f(v)\,|\le L\,|u-v|.
\]
Equivalently,
\[
f(u)\le f(v)+\nabla f(v)(u-v)+\frac{L}{2}(u-v)^2 .
\]
\end{assumption}
\begin{assumption}[Bounded Gradients]\label{as:bound}
There exists $G_{\max}>0$ such that $|g_k|=|\nabla f(w_k)|\le G_{\max}$ for all $k$.
\end{assumption}
\begin{assumption}[Bounded Domain]\label{as:diam}
The distance between any two iterates is bounded:
\[
|w-w'|\le D,\qquad\forall w,w'\in\{w_0,\dots,w_T\},
\]
i.e. the feasible set has diameter $D$.
\end{assumption}

All constants $\alpha,\varepsilon,G_{\max},L,D$ are assumed known and
strictly positive.

\section*{Main Convergence Theorem}
\begin{theorem}[AdaGrad convergence for convex smooth $f$]\label{thm:main}
Let $\{w_k\}_{k=0}^{T}$ be generated by the scalar AdaGrad algorithm with
$\varepsilon>0$.
Define the averaged iterate $\bar w_T:=\frac{1}{T}\sum_{k=1}^{T} w_k$.
Under Assumptions \ref{as:convex}–\ref{as:diam}, the following bound holds
\[
\boxed{
f(\bar w_T)-f(w^{\star})
\;\le\;
\frac{D\,G_{\max}}{\alpha}\,
\frac{1}{\sqrt{T}} \;+\;
\frac{\alpha}{2\sqrt{T}}\,
\Bigl(1+\frac{G_{\max}^{2}}{\varepsilon}\Bigr)
}
\]
where $w^{\star}\in\arg\min f$.
Consequently,
\[
f(\bar w_T)-f(w^{\star}) = O\!\Bigl(\frac{1}{\sqrt{T}}\Bigr).
\]
\end{theorem}

\section*{Theoretical Properties}
\begin{enumerate}[label=\arabic*.]
\item \textbf{Stability.} The denominator $\sqrt{G_k+\varepsilon}$ is non‑decreasing,
hence $\eta_k$ is non‑increasing. This prevents exploding updates.
\item \textbf{Hyper‑parameter sensitivity.}
\begin{itemize}
\item Larger $\alpha$ accelerates early progress but may increase the constant
term $\frac{\alpha}{2\sqrt{T}}(1+G_{\max}^{2}/\varepsilon)$.
\item $\varepsilon$ only affects the very first steps; setting $\varepsilon\ll G_{\max}^{2}$
does not change the asymptotic rate.
\end{itemize}
\item \textbf{Comparison to SGD.}
For constant stepsize $\eta$, SGD obtains $O(1/\sqrt{T})$ only under a
bounded variance assumption. AdaGrad achieves the same rate without
any stochasticity, and the bound automatically adapts to the observed
gradient magnitudes.
\item \textbf{Special case $G_{\max}=0$.}
If all gradients vanish after some iteration, $G_k$ stops growing,
$\eta_k$ stabilises, and the algorithm reaches the exact optimum in a
finite number of steps.
\end{enumerate}

\section*{Discussion}
The proof mirrors the classic AdaGrad analysis for online convex optimisation
but specialised to the deterministic setting. The key idea is to bound the
inner product $\langle g_k,w_k-w^{\star}\rangle$ by a telescoping sum that
involves the adaptive stepsizes. The smoothness assumption is only needed
to relate function values to gradients in the final averaging step.
Because the stepsizes shrink as $1/\sqrt{\sum_{i=1}^{k}g_i^{2}}$, the bound
depends on the \emph{empirical} magnitude of the gradients, yielding
potentially tighter constants than a worst‑case SGD analysis.

\newpage
\section*{Preliminaries (Definitions and Lemmas)}
\begin{lemma}[Quadratic expansion]\label{lem:quad}
For any $a,b\in\R$ and any $\eta>0$,
\[
(a-\eta b)^{2}=a^{2}-2\eta a b+\eta^{2}b^{2}.
\]
\end{lemma}
\begin{lemma}[Cauchy–Schwarz for scalars]\label{lem:cs}
For any $x,y\in\R$,
\[
|xy|\le\frac{1}{2}\bigl(x^{2}+y^{2}\bigr).
\]
\end{lemma}
\begin{lemma}[Bound on accumulated stepsizes]\label{lem:sumeta}
For the AdaGrad stepsizes $\eta_k=\alpha/\sqrt{G_k+\varepsilon}$,
\[
\sum_{k=1}^{T}\eta_k\, g_k^{2}
\;=\;
\alpha\sum_{k=1}^{T}\frac{g_k^{2}}{\sqrt{G_k+\varepsilon}}
\;\le\;
2\alpha\bigl(\sqrt{G_T+\varepsilon}-\sqrt{\varepsilon}\bigr).
\]
\end{lemma}
\begin{proof}
Define $h_k:=\sqrt{G_k+\varepsilon}$. 
Since $G_k=G_{k-1}+g_k^{2}$,
\[
h_k-h_{k-1}
=\frac{G_k+\varepsilon-(G_{k-1}+\varepsilon)}{h_k+h_{k-1}}
=\frac{g_k^{2}}{h_k+h_{k-1}}
\ge\frac{g_k^{2}}{2h_k}.
\]
Rearranging gives $g_k^{2}\le 2h_k(h_k-h_{k-1})$.
Multiplying by $\alpha/h_k$ and summing yields
\[
\sum_{k=1}^{T}\eta_k g_k^{2}
\le 2\alpha\sum_{k=1}^{T}(h_k-h_{k-1})
=2\alpha\bigl(h_T-h_0\bigr)
=2\alpha\bigl(\sqrt{G_T+\varepsilon}-\sqrt{\varepsilon}\bigr).
\]
\end{proof}

\section*{Main Proof (Step‑by‑Step Derivation)}
We prove Theorem \ref{thm:main}. Throughout the proof we denote the optimal
point by $w^{\star}\in\arg\min f$.

\bigskip
\noindent\textbf{[Step 1]} (Start from the squared distance recursion.)  
Using Lemma \ref{lem:quad} with $a=w_k-w^{\star}$, $b=g_k$, and
$\eta=\eta_k$,
\begin{align}
\bigl(w_{k+1}-w^{\star}\bigr)^{2}
&= \bigl(w_k-w^{\star}-\eta_k g_k\bigr)^{2}\nonumber\\
&= (w_k-w^{\star})^{2}
-2\eta_k g_k (w_k-w^{\star})
+\eta_k^{2} g_k^{2}. \label{eq:dist}
\end{align}

\noindent\textbf{[Step 2]} (Rearrange to isolate the inner product.)  
Move the first term to the left and divide by $2\eta_k$:
\begin{align}
g_k (w_k-w^{\star})
&= \frac{(w_k-w^{\star})^{2}-(w_{k+1}-w^{\star})^{2}}{2\eta_k}
+\frac{\eta_k}{2}\, g_k^{2}. \label{eq:inner}
\end{align}

\noindent\textbf{[Step 3]} (Use convexity.)  
By convexity (Assumption \ref{as:convex}),
\[
f(w_k)-f(w^{\star})\le g_k (w_k-w^{\star}).
\]
Insert \eqref{eq:inner}:
\begin{align}
f(w_k)-f(w^{\star})
&\le
\frac{(w_k-w^{\star})^{2}-(w_{k+1}-w^{\star})^{2}}{2\eta_k}
+\frac{\eta_k}{2}\, g_k^{2}. \label{eq:conv}
\end{align}

\noindent\textbf{[Step 4]} (Sum over $k=0,\dots,T-1$.)  
Summation telescopes the first fraction:
\begin{align}
\sum_{k=0}^{T-1}\!\bigl[f(w_k)-f(w^{\star})\bigr]
&\le
\frac{(w_0-w^{\star})^{2}}{2\eta_0}
+\sum_{k=1}^{T-1}\!\Bigl(\frac{1}{2\eta_k}-\frac{1}{2\eta_{k-1}}\Bigr)(w_k-w^{\star})^{2}
+\frac12\sum_{k=0}^{T-1}\eta_k g_k^{2}. \label{eq:sum}
\end{align}
Since $\eta_k$ is non‑increasing, the coefficient
$\frac{1}{2\eta_k}-\frac{1}{2\eta_{k-1}}\le 0$, thus we can drop the middle term
to obtain an upper bound:
\begin{align}
\sum_{k=0}^{T-1}\!\bigl[f(w_k)-f(w^{\star})\bigr]
&\le
\frac{(w_0-w^{\star})^{2}}{2\eta_0}
+\frac12\sum_{k=0}^{T-1}\eta_k g_k^{2}. \label{eq:bound1}
\end{align}

\noindent\textbf{[Step 5]} (Bound the initial denominator.)  
Because $\eta_0=\alpha/\sqrt{g_0^{2}+\varepsilon}\ge \alpha/\sqrt{G_{\max}^{2}+\varepsilon}$,
\[
\frac{1}{\eta_0}\le \frac{\sqrt{G_{\max}^{2}+\varepsilon}}{\alpha}
\le \frac{G_{\max}+\sqrt{\varepsilon}}{\alpha}.
\]

\noindent\textbf{[Step 6]} (Bound the accumulated term using Lemma \ref{lem:sumeta}).  
From Lemma \ref{lem:sumeta},
\[
\sum_{k=0}^{T-1}\eta_k g_k^{2}
\le
2\alpha\bigl(\sqrt{G_T+\varepsilon}-\sqrt{\varepsilon}\bigr)
\le
2\alpha\bigl(\sqrt{T}\,G_{\max}+\sqrt{\varepsilon}\bigr),
\]
where we used $G_T=\sum_{i=0}^{T-1}g_i^{2}\le T G_{\max}^{2}$.

\noindent\textbf{[Step 7]} (Insert the bounds into \eqref{eq:bound1}.)  
Denote $D:=\max_{k}|w_k-w^{\star}|\le |w_0-w^{\star}|+ \sum_{i=0}^{k-1}\eta_i|g_i|
\le D$ by Assumption \ref{as:diam}. Then
\begin{align}
\sum_{k=0}^{T-1}\!\bigl[f(w_k)-f(w^{\star})\bigr]
&\le
\frac{D^{2}}{2}\,\frac{G_{\max}+\sqrt{\varepsilon}}{\alpha}
+\alpha\bigl(\sqrt{T}\,G_{\max}+\sqrt{\varepsilon}\bigr). \label{eq:sumfinal}
\end{align}

\noindent\textbf{[Step 8]} (Convert the sum bound to a bound on the average iterate.)  
Because $f$ is convex,
\[
f(\bar w_T)-f(w^{\star})
\;\le\;
\frac{1}{T}\sum_{k=0}^{T-1}\!\bigl[f(w_k)-f(w^{\star})\bigr].
\]
Dividing \eqref{eq:sumfinal} by $T$ yields
\begin{align}
f(\bar w_T)-f(w^{\star})
&\le
\frac{D^{2}}{2\alpha}\,\frac{G_{\max}+\sqrt{\varepsilon}}{T}
+\frac{\alpha G_{\max}}{\sqrt{T}}
+\frac{\alpha\sqrt{\varepsilon}}{T}. \label{eq:avg}
\end{align}

\noindent\textbf{[Step 9]} (Simplify constants.)  
Dropping the $1/T$ terms (which are dominated by $1/\sqrt{T}$ for $T\ge1$) we obtain
\[
f(\bar w_T)-f(w^{\star})
\le
\frac{D\,G_{\max}}{\alpha}\,\frac{1}{\sqrt{T}}
+\frac{\alpha}{2\sqrt{T}}\Bigl(1+\frac{G_{\max}^{2}}{\varepsilon}\Bigr),
\]
which is exactly the bound stated in Theorem \ref{thm:main}. ∎

\section*{Rate Derivation (With Constants)}
The bound can be written as
\[
f(\bar w_T)-f(w^{\star})
\le
\underbrace{\frac{D\,G_{\max}}{\alpha}}_{C_1}\frac{1}{\sqrt{T}}
\;+\;
\underbrace{\frac{\alpha}{2}\Bigl(1+\frac{G_{\max}^{2}}{\varepsilon}\Bigr)}_{C_2}\frac{1}{\sqrt{T}}
\;=\;
\frac{C_1+C_2}{\sqrt{T}}.
\]
Thus the asymptotic rate is $O(1/\sqrt{T})$ with an explicit constant
$C_1+C_2$ that depends linearly on the diameter $D$ and inversely on the
initial stepsize $\alpha$, while $C_2$ captures the influence of the
regulariser $\varepsilon$.

\section*{Special Cases}
\begin{enumerate}[label=\alph*)]
\item \textbf{Exact gradients vanish after $k^{\star}$.}  
If $g_k=0$ for all $k\ge k^{\star}$, then $G_k$ stops growing, $\eta_k$ becomes
constant for $k\ge k^{\star}$, and the algorithm reaches $w^{\star}$ in at most
$k^{\star}+1$ steps (the bound becomes zero).

\item \textbf{Choosing $\alpha = D G_{\max}$}.  
The two terms in the bound become equal, giving
\[
f(\bar w_T)-f(w^{\star})\le
\frac{D\,G_{\max}}{D\,G_{\max}}\frac{1}{\sqrt{T}}
+\frac{D\,G_{\max}}{2\sqrt{T}}\Bigl(1+\frac{G_{\max}^{2}}{\varepsilon}\Bigr)
= O\!\Bigl(\frac{1}{\sqrt{T}}\Bigr).
\]

\item \textbf{Limit $\varepsilon\to 0$.}  
The bound reduces to
\[
f(\bar w_T)-f(w^{\star})\le
\frac{D\,G_{\max}}{\alpha}\frac{1}{\sqrt{T}}
+\frac{\alpha G_{\max}^{2}}{2\sqrt{T}},
\]
showing the regulariser only affects the constant term and not the asymptotic order.
\end{enumerate}

\end{document}